\documentclass[11pt]{article}

\usepackage[utf8]{inputenc}
\usepackage[margin=1in]{geometry}
\usepackage{amsmath}
\usepackage{amssymb}
\usepackage{physics}
\usepackage{xcolor}
\usepackage{mdframed}
\usepackage{hyperref}
\hypersetup{colorlinks=true, allcolors=blue}
\usepackage{enumitem}

% Reviewer comment: boxed italic block
\newmdenv[
  linecolor=gray!55,
  linewidth=0.6pt,
  backgroundcolor=gray!6,
  leftmargin=0pt, rightmargin=0pt,
  innertopmargin=6pt, innerbottommargin=6pt,
  skipabove=8pt, skipbelow=4pt
]{revcomment}

\newcommand{\comment}[1]{\begin{revcomment}\itshape #1\end{revcomment}}
\newcommand{\response}[1]{\noindent\textbf{Response.}\ #1\par\vspace{6pt}}

\setlength{\parskip}{4pt}
\setlength{\parindent}{0pt}

\title{\vspace{-2.5em}Response to Reviewers\\[2pt]
\large Manuscript: ``Non-Equilibrium Dynamics of the Time-Dependent Excitonic Coupling in Fluorescent Protein Dimers''}
\author{R.\ Christie, C.\ Murray, Y.\ Kim, J.\ Joo}
\date{}

\begin{document}
\maketitle
\vspace{-2.5em}

\section*{Cover note to the Editor}

We thank the Editor and the four reviewers for their careful and constructive
assessment, and we are grateful for the invitation to transfer the manuscript to
\emph{The Journal of Physical Chemistry B}. We have revised the manuscript
substantially in response to the reports. The principal changes are:

\begin{itemize}[leftmargin=1.4em, itemsep=2pt]
  \item \textbf{Higher-level electronic structure.} We have strengthened the
  electronic-structure treatment underlying the coupling by computing the site energy
  and transition density with a higher-order correlated wavefunction method,
  DLPNO-STEOM-CCSD. Because the anionic chromophore carries appreciable
  double-excitation character, this correlated treatment captures physics beyond a
  single-excitation description and brings the predicted site energy into closer
  agreement with the measured absorption. The STEOM assignment is validated on the
  tractable gas-phase chromophore against an explicit triples correction
  [EOM-CCSD(fT)] and ADC(2).
  \item \textbf{Conformational sampling.} The coupling is now characterised not at a
  single minimised geometry but as a thermal distribution over a
  1000-frame unrestrained NVT ensemble of the tandem dimer.
  \item \textbf{Multipole origin of the PDA breakdown.} We add an explicit
  multipole decomposition (dipole through octupole) quantifying which terms contribute
  and demonstrating that the enhancement resides in the short-range Coulomb coupling
  between the extended (non-overlapping) transition densities.
  \item \textbf{A spectroscopic observable.} We connect the coupling and the
  open-system parameters to an absorption/CD lineshape (new Figure~5), yielding the bisignate CD
  couplet characteristic of the dimer.
  \item \textbf{Methodological transparency.} Grid spacing and QM region sizes are now
  stated explicitly, all software is cited, and the entire workflow is reproduced by a
  single released orchestrator, \texttt{reproduce\_paper.py}.
\end{itemize}

Point-by-point responses to each reviewer follow, each beginning on a new page so
that they may be forwarded individually. Reviewer comments are reproduced in the
shaded boxes.

\newpage
\section*{Response to Reviewer 1}

We thank the reviewer for recognising the two central contributions (the
timescale-separation argument and the near-field enhancement of the coupling) and for
the constructive suggestions, which we address in turn.

\comment{I did not see much meaningful integration between the open quantum systems
simulations and the quantum chemistry part\ldots\ the approach of generating dynamics
is not used to, for example, calculate a spectroscopic signature such as an absorption
(or emission) spectrum.}

\response{We agree, and we have added exactly this link. From the computed coupling,
its conformational spread (which fixes the inhomogeneous Gaussian width) and the
homogeneous dephasing width $1/(2\pi c\,T_2^{*})$, we now construct the excitonic
absorption and circular-dichroism lineshapes, modelling each exciton band as a Voigt
profile (new Figure~5; tool \texttt{absorption\_cd\_spectra.py}). This produces the bisignate CD couplet
characteristic of an excitonic dimer, with an apparent splitting ($2|J|\approx 221~\mathrm{cm^{-1}}$) of the same order as, and just below, the
experimentally reported window ($262$--$372~\mathrm{cm^{-1}}$), so the open-system parameters are now tied to a
measurable observable rather than being purely illustrative. The dynamics thus feed
directly into the spectroscopic signature and share the same coupling and dephasing
inputs.}

\comment{To what extent are the findings and claims specific to the protein dimer
under consideration?\ \ldots\ Wouldn't at least the (only slowly kicking in) solvent
shielding be expected in other aqueous complexes\ldots?}

\response{The timescale-separation argument is not specific to Venus: aqueous Debye
relaxation universally lies in the picosecond range, while optical absorption and
dephasing are far faster, so the optical-limit imprinting of the coupling is expected
in any water-exposed multichromophore complex. What the fluorescent-protein scaffold adds is a
reduction of the fluctuation \emph{amplitude}: the covalently tethered construct holds the
A206 interface docked so that the coupling varies by ${\sim}3\%$ across the thermal
ensemble, not a change in the underlying separation of timescales.}

\comment{It is intriguing that the point dipole approximation seems to be out by a
factor of five\ldots\ are broader lessons to be learned regarding when we should
generally be more careful in applying the point dipole approximation?}

\response{Yes. We now decompose the coupling into a multipole
series (dipole, quadrupole, octupole and the corresponding cross terms). The
dipole--dipole term reproduces the PDA exactly, but the series converges only slowly:
through octupole it recovers only ${\sim}28\%$ of the full coupling. The lesson,
now stated explicitly, is that when the transition density is delocalised and its
nearest inter-chromophore approach is comparable to the molecular size, even a
multipole-corrected dipole model is inadequate and the full density integral is
required, a criterion that generalises beyond this system.}

\comment{Regarding the solvent polarisation (Debye) time, aren't there also fast
timescales leading to a partial polarisation?}

\response{Indeed. We now note in the Summary that the single-exponential Debye form
omits the fast librational and intramolecular components of the water response, which
act on tens of femtoseconds and provide a partial polarisation before full nuclear
reorganisation. Consequently the optical limit is best regarded as an upper bound on
the early-time coupling; it does not alter the central conclusion, since these fast
components still lie far below the picosecond Debye time.}

\comment{Given the relatively tight packing\ldots\ the solvent distribution will be
quite asymmetric. Is it right to assume this affects the strict validity of Eqs 3, 4,
or renders them into an approximation?}

\response{The reviewer is correct: the isotropic-continuum screening in the point-dipole coupling and the
effective dielectric-relaxation model is an approximation, and the asymmetric, partially excluded solvent
distribution in the inter-barrel gap means a position-dependent or atomistic,
explicitly polarisable embedding would be more rigorous. We now state this as a
limitation and point to such embeddings as the appropriate refinement. Importantly,
however, the quantity imprinted at absorption, $J(0)$, is set by the optical dielectric
$\varepsilon_{\rm opt}$ alone; sweeping the static permittivity $\varepsilon_{\rm eq}$
over $4$--$78$ leaves $J(0)$ invariant and rescales only the long-time relaxed limit,
so the asymmetry bears on the late-time damping rather than on the Davydov splitting.}

\comment{I was not sure how valuable ``Algorithm 1'' is\ldots\ and what exactly one
should take away from Table 1.}

\response{We have reframed both. Algorithm 1 is now presented purely as a compact,
reproducible map of the released pipeline and may be moved to the Supporting
Information at the Editor's discretion. Table 1 has been replaced by a concise
cross-method comparison of the bright-state energy (TDDFT vs.\ STEOM-CCSD vs.\
experiment).}

\newpage
\section*{Response to Reviewer 2}

We thank the reviewer for the positive assessment of the timescale proposition and for
the recommendation to reconsider the work as a full \emph{Journal of Physical
Chemistry} article, which we have accepted.

\comment{The lack of comparison between these dynamics and experimental observables
means it is hard to distinguish whether the model is qualitatively correct or just one
of many possible toy-models of dephasing.}

\response{We have added the missing comparison. The coupling and dephasing now feed a
Voigt absorption/CD lineshape (new Figure~5) whose bisignate couplet can be compared with the observed
CD signature. In addition, we clarify that $T_2^{*}=60$\,fs is an order-of-magnitude
estimate transferred from analogous systems, and we show by sweeping $T_2^{*}$ over
$20$--$200$\,fs that the central conclusion is insensitive to its precise value: all
resulting coherence-decay times remain far below the picosecond Debye time.}

\comment{Contemporary works regularly use molecular dynamics ensembles rather than a
single crystal structure\ldots\ (e.g., DOI: 10.1021/acs.jpcc.7b00833). In addition,
the environment is often modelled more accurately, for instance by including its
polarizability (e.g., DOI: 10.3389/fmolb.2023.1268278).}

\response{We have adopted the ensemble approach the reviewer highlights. The coupling
is now averaged over a 1000-frame unrestrained NVT trajectory of the tandem dimer, giving
$J = 111 \pm 3~\mathrm{cm^{-1}}$ (a $\sim3\%$ spread lower bound at 300K) rather than a single-structure
value. On accuracy, the transition density itself is now taken from a correlated
wavefunction method (DLPNO-STEOM-CCSD) rather than TDDFT. We note that the environment's fast electronic polarisability is already represented
at the continuum level, through the optical dielectric $\varepsilon_{\rm opt}$ that screens the
coupling; an explicitly polarisable embedding would add the explicit, position-dependent
atomistic response rather than electronic polarisability per se, and would further refine
the screening at this separation. We now
cite such approaches (including the frontiers reference) and identify polarisable
embedding as the appropriate next refinement, while noting that the absorption-time
coupling $J(0)$ is governed by $\varepsilon_{\rm opt}$ and is invariant to the static
permittivity.}

\comment{Intermediate excitonic coupling is stated to be ``at odds with the biological
environment inducing strong decoherence'', but values such as the present one, and even
larger, are common in photosynthetic pigment--protein complexes (e.g.,
10.1146/annurev.physchem.54\ldots).}

\response{Good point. We agree that couplings of this magnitude, and
larger, are well documented in photosynthetic pigment--protein complexes, and the
revised Introduction now states this explicitly and cites it, rather than presenting the
magnitude of $J$ as itself surprising. We have correspondingly reframed the open
question as the room-temperature \emph{persistence} of the coherent spectroscopic
signatures, which is what the separation-of-timescales argument addresses. We are
grateful to the reviewer for prompting this clarification.}

\comment{The authors present as a result that the point-dipole approximation is not
accurate for distances comparable to the molecular size, which is nowadays a well
established fact (e.g., 10.1021/acs.jpcb.7b09535).}

\response{We no longer present the breakdown itself as novel. The Introduction
and Summary now cite this established result, and we recast our contribution as (i)
\emph{quantifying} the magnitude of the effect for this fluorescent-protein dimer and
(ii) identifying its multipolar origin through an explicit decomposition, which shows
that the enhancement lies in the short-range Coulomb coupling between the extended
(non-overlapping) transition densities rather than in a few leading multipoles.}

\comment{My only suggestion is to add small arrows in Figure 1a to indicate the time
progression of the two trajectories.}

\response{We have clarified the time progression in the caption of Figure~1: both paths
begin at the bright state $\ket{+}$ on the equator, the individual (red) trajectory
executes surface-preserving stochastic hops, and the ensemble average (blue) contracts
radially inward toward the maximally mixed centre; panels (b) and (c) run left to right
in time. We attempted to add arrowheads through the stochastic trajectory but the result in our opinion looked a bit blob like.}

\newpage
\section*{Response to Reviewer 3}

We thank the reviewer for the high assessment of significance and novelty and for the
specific, actionable concerns, which we have addressed.

\comment{The manuscript provides only a single static calculation based on one
minimized geometry. No conformational sampling, molecular dynamics averaging, or
analysis of coupling fluctuations is performed.}

\response{We have added conformational sampling. The coupling is now evaluated on 1000
snapshots from an unrestrained NVT trajectory of the tandem dimer, giving
$J = 111 \pm 3~\mathrm{cm^{-1}}$ with the centroid separation tightly clustered so the reported value is
representative of the thermal ensemble rather than one configuration among many.}

\comment{The manuscript repeatedly discusses environmental fluctuations and
decoherence, yet the coupling calculations completely neglect structural dynamics.}

\response{This is now remedied by the ensemble calculation above, which explicitly
propagates structural fluctuations into the coupling and reports their statistics.}

\comment{The use of bulk water parameters ($\varepsilon = 78$ and $\tau_D = 8.3$\,ps)
may be inappropriate for chromophores embedded inside a protein $\beta$-barrel. Is this
approximation justified?}

\response{We now discuss this explicitly as a limitation. The barrel-enclosed
chromophore is only partially solvent-exposed, so the effective static permittivity is
lower and heterogeneous. Crucially, though, the absorption-time coupling $J(0)$ depends
only on the optical dielectric $\varepsilon_{\rm opt}$; a sweep of $\varepsilon_{\rm eq}$
over $4$--$78$ leaves $J(0)$ invariant and affects only the long-time relaxed limit.
The choice of static permittivity therefore does not affect the Davydov splitting set
at absorption.}

\comment{A dephasing time of 60 fs is adopted without rigorous justification\ldots\
There is no evidence that the same dephasing rate applies to the Venus dimer.}

\response{We now state plainly that no $T_2^{*}$ has been measured for the Venus dimer
and that $60$\,fs is an order-of-magnitude estimate transferred from analogous solvated
pigment--protein systems. It enters only as an upper bound on the dephasing timescale.
A sweep over $T_2^{*} = 20$--$200$\,fs shows that the separation of timescales, and
hence the central claim, is preserved throughout, so the argument does not rely on the
precise value.}

\comment{Although Figures 2 and 3 are attractive, they provide limited quantitative
information and would benefit from additional quantitative analysis.}

\response{We have added two quantitative analyses grounded in these densities: the
Cartesian multipole decomposition of the coupling (dipole/quadrupole/octupole
contributions and their slow convergence) and the absorption/CD lineshape derived from
the coupling (new Figure~5). Together these convert the qualitative density pictures into quantitative
statements about the coupling and its spectroscopic signature.}

\comment{``rigorous QM rationale,'' ``fundamentally reconciled,'' and ``robust
coupling,'' could be toned down.}

\response{We have moderated the language throughout. The overstated phrases have been
removed or rephrased, and the remaining uses of ``robust'' are now quantitatively
justified.}

\newpage
\section*{Response to Reviewer 4}

We thank the reviewer for the detailed technical reading. Each point has been
addressed, and several prompted substantive improvements to the electronic structure
and the numerical documentation.

\comment{The time-dependent permittivity in Eq. 4 adopts an equilibrium dielectric
constant of 78 for liquid water, which is not relevant for the protein environment.}

\response{We now discuss this limitation explicitly and show that the central result is
insensitive to it: the coupling imprinted at absorption, $J(0)$, is set by the optical
dielectric alone, and varying $\varepsilon_{\rm eq}$ from $4$ (protein interior) to
$78$ (bulk water) leaves $J(0)$ unchanged, affecting only the long-time relaxed limit.}

\comment{The authors mentioned ``the rapid, smooth decay of off-diagonal coherence''.
But it cannot be clearly seen in Figure 1b.}

\response{We have added a second row of panels to Figure~1; in particular,
\textbf{Figure~1(e)} now shows explicitly the decay of the off-diagonal density-matrix
element (the ensemble-averaged site-basis coherence $|\rho_{12}|$), and the caption points
to its rapid sub-$100$\,fs decay in panels (b) and (e). The surrounding text gives the
decay time from the $T_2^{*}$ sweep, so the off-diagonal decay is now clearly identifiable
in the figure.}

\comment{They should at least compute the monomer quadrupole and octupole moments of
the transition density, and show how dipole--quadrupole, quadrupole--quadrupole, and
higher-order couplings contribute.}

\response{We have now done exactly this. A new analysis computes the transition dipole,
quadrupole and octupole moments and the full set of inter-site interaction terms
(dipole--dipole, dipole--quadrupole, quadrupole--dipole, dipole--octupole,
quadrupole--quadrupole, octupole--dipole). Summing through octupole recovers only
${\sim}28\%$ of the full coupling (dipole alone ${\sim}23\%$, through quadrupole
${\sim}27\%$), demonstrating that the interaction is not captured by any low-order
multipole truncation. We emphasise that this is a purely Coulombic (transition-density)
effect: the two densities do not overlap (their facing tails approach to ${\sim}9$~\AA,
the nuclei to ${\sim}15$~\AA), so it is distinct from short-range Dexter exchange, which
remains negligible at these separations.}

\comment{For the computation of transition density coupling\ldots\ the spacing between
grid points should be specified. As mentioned in Ref. 13, one can compute such
couplings analytically\ldots\ This avoids the use of a grid representation.}

\response{We now specify the grid: a uniform $200\times200\times200$ Cartesian grid
with spacing $\approx 0.19~\mathrm{\AA}$ (voxel volume $\approx 7\times10^{-3}~
\mathrm{\AA^{3}}$), of which $\sim3\times10^{5}$ significant voxels enter the Coulomb
double sum. We also acknowledge in the Methods that these couplings can be evaluated
analytically from atom-centred transition charges, and note that our dense-grid
integral converges to the same value while additionally supplying the volumetric
densities shown in the figures.}

\comment{For QM/MM modeling, one should specify the number of QM atoms.}

\response{The QM region is now stated explicitly: both the single-excitation TDDFT
reference and the correlated DLPNO-STEOM-CCSD calculation use the same 44-atom region
(the deprotonated CR2 chromophore plus the $\pi$-stacked Tyr203, including hydrogen
link-atom caps), embedded in the identical MM point-charge field, so the two methods
are compared on an identical structure and embedding.}

\comment{It is not sufficient to compare the values of the couplings for a single
frame\ldots\ the authors should perform constant-temperature simulations to collect a
set of system configurations, and analyze the couplings for those configurations.}

\response{We agree and have implemented this. The coupling is now analysed over a
1000-frame unrestrained NVT ensemble of the tandem dimer rather than a single minimised structure, yielding
$J = 111 \pm 3~\mathrm{cm^{-1}}$ and demonstrating that the value is stable
($\sim3\%$ variation) with respect to thermal geometric disorder.}

\end{document}
